Denne oppgaven handler om stokastiske impulskontrollproblemer hor hoppeprossesser over en endelig tidshorisont. Teoretiske løsningsmåter som dynamisk programmering og den kvasi-variasjonsulikheten (QVI) er gjennomgått. Fordi løsninger av QVI-er generelt sett ikke er glatte, brukes viskositetsløsninger for å bevise at løsningen på QVI-en er løsningen på impulskontrollproblemet. QVI-en løses numerisk ved bruk av differansemetoder. 

Impulsekontrollproblemet kobles sammen med dyp forsterkningslæring (DRL) ved hjelp av dynamic programmering. To simulasjonsbaserte DRL algoritmer er tatt i betraktning: den verdibaserte Double Deep Q-Network-metoden (DDQN) og den polise-baserte Proximal Policy Optimisation-algoritmen som er utvidet med selvimitasjonslæring (PPO-SIL). 

Til slutt anvendes metodene på to problemer innen finans. Det første problemet handler om den optimale kontrollem av valutakurset med kurtasje, , hvor både DDQN- og PPO-SIL-algoritmene sammenlignes med løsningen fra endelig-differansemetoden. I det andre problemet, som er en dynamisk porteføljeallokeringsproblem med kurtasje, brukes PPO-SIL til å forvalte en portefølje med 20 aksjer over en tidsperiode av ett år. Aksjepriser antas til å være korrelerte, multidimensjonale varianse-gamma-prosesser med en felles subordinator, og Europeiske kjøpsopsjoner brukes til å estimere parameterne inni aksjeprismodellen. PPO-modellen oppnår en gjennomsnittlig kumulativ avkastning på rundt 10\% og en gjennomsnittlig årlig Sharpe-ratio på 3.26. 